\begin{abstract}
Although large-scale unsupervised language models (LMs) acquire extensive world knowledge and certain reasoning capabilities, exerting fine-grained control over their outputs remains challenging due to the fully unsupervised nature of their training. Approaches designed to enhance steerability typically rely on gathering human judgments about the relative merits of model outputs, followed by adapting the unsupervised LM through fine-tuning to reflect these preferences, often leveraging reinforcement learning from human feedback (RLHF). Yet, RLHF involves a multi-stage and frequently unstable process, beginning with the training of a reward model that encapsulates human preferences and subsequently using reinforcement learning to update the LM so as to optimize the inferred reward while maintaining proximity to the original parameters. In this work, we propose a novel reward model parameterization for RLHF that allows for a closed-form derivation of the optimal policy, thereby reducing the RLHF problem to the minimization of a standard classification loss. This leads to our method, termed \textit{Direct Preference Optimization} (DPO), which is characterized by its robustness, strong empirical performance, and low computational cost. DPO removes the necessity for sampling from the LM during the adaptation stage and requires minimal hyperparameter adjustment. Experimental results demonstrate that DPO fine-tunes LMs to match or surpass the performance of established preference alignment techniques. In particular, DPO outperforms RLHF with PPO in regulating generative sentiment, while achieving comparable or superior results in summarization and single-turn dialogue tasks, all within a framework that is substantially easier to implement and train.
\end{abstract}

\section{Introduction}
Large-scale unsupervised language models (LMs), when exposed to massive corpora, often develop unexpected proficiencies~\citep{chowdhery2022palm, brown2020language, touvron2023llama,bubeck2023sparks}. Yet, the data used in their training originates from human authors with diverse intentions, backgrounds, and competencies. Not every objective or expertise captured in the dataset is necessarily desirable for downstream applications. For instance, while an AI tool aimed at assisting with code should ideally \textit{recognize} typical programming errors to aid in fixing them, it should preferentially \textit{emulate} the infrequent but superior programming skill present within its training corpus when producing code. By the same token, although it is beneficial for a language model to possess \textit{knowledge} of a widespread misconception held by half the population, we would prefer that the model does not propagate this misconception as fact in 50\% of its responses regarding the topic. Thus, determining the \emph{intended outputs and behavioral patterns} a model should express—drawn from its broad \textit{repository of knowledge and capabilities}—is a fundamental step toward developing AI that is safe, effective, and amenable to control~\citep{ouyang2022training}. While most current strategies employ reinforcement learning (RL) to tune language models according to human preferences, we will demonstrate that the objective function underlying these RL approaches can, in fact, be exactly minimized using a straightforward binary cross-entropy loss, streamlining the process of learning from human preferences.

Broadly, the dominant practice for imparting appropriate behavioral biases to language models is to utilize carefully assembled datasets of human preferences, designed to capture human judgments of safety and helpfulness. This phase of preference alignment is typically carried out subsequent to a comprehensive, unsupervised pre-training stage over extensive text collections. Supervised fine-tuning on examples of exemplary human responses is the most direct technique, but reinforcement learning from human (or artificial) feedback (RLHF/RLAIF; \citep{christiano2017deep,bai2022constitutional}) has emerged as the most influential family of methods. RLHF entails training a reward model on human-preference data, followed by RL optimization of the LM’s output policy, ensuring it yields high-reward answers without diverging markedly from the initial, pre-trained model. While this approach produces models notable for their dialogue and programming prowess, RLHF pipelines are substantially more intricate than those of supervised learning, necessitating the training of several LMs and the integration of online sampling from the model policy during training—ultimately resulting in notable computational overhead.

This work introduces a method for optimizing language models according to human preferences without relying on explicit reward modeling or reinforcement learning techniques. We introduce \textit{{Direct Preference Optimization} (DPO)}, an approach that implicitly targets the same objective as standard RLHF algorithms—namely, reward maximization subject to a KL-divergence constraint—while remaining easy to implement and efficient to train. At a high level, DPO works by amplifying the log-likelihood ratio of preferred over non-preferred outputs, yet it also employs an adaptive, per-sample importance weighting mechanism. This dynamic weighting mitigates the mode collapse observed with straightforward application of probability ratio objectives. DPO, like prior techniques, makes use of a theoretical preference model—such as the Bradley-Terry model \cite{bradley1952rankanalysis}—for quantifying the alignment between a candidate reward function and observed preference data. In contrast to approaches that use this model to generate a preference loss for reward model training followed by policy optimization, DPO leverages a variable transformation to express the preference loss directly in terms of the policy. When provided with a collection of human-annotated preference data for model outputs, DPO directly trains the policy using a standard binary cross-entropy loss, effectively yielding the optimal policy for an implicit reward function inferred from these preferences.

Our primary contribution is the development of Direct Preference Optimization (DPO), an algorithm that allows for preference-driven language model training without the need for reinforcement learning. Through empirical evaluation, we demonstrate that DPO achieves performance on par with, or better than, established methods—including PPO-based RLHF—across preference-sensitive tasks such as sentiment adjustment, text summarization, and conversational modeling, utilizing language models scaling up to 6B parameters.

\section{Related Work}

Large-scale self-supervised language models are capable of performing certain tasks in a zero-shot manner \citep{radford2019language}, or when provided with few-shot prompts \citep{gpt3,megatron,chowdhery2022palm}. Nonetheless, these models demonstrate enhanced effectiveness on downstream applications and better alignment with user expectations when they undergo fine-tuning using collections of instructions and human-generated responses \citep{mishra-etal-2022-cross,sanh2022multitask,chung2022scaling,thoppilan2022lamda}. This process, known as ‘instruction-tuning,’ extends LLMs’ ability to generalize beyond the set of instructions present in the training data, thereby making them broadly more useful \citep{chung2022scaling}. Although instruction tuning has proven to be highly effective, obtaining \textit{relative} human assessments of response quality tends to be less demanding than collecting expert demonstrations. Consequently, later research has focused on fine-tuning LLMs using datasets composed of human preference judgments, which has been shown to improve task proficiency across domains including translation \citep{kreutzer-etal-2018-reliability}, summarization \citep{stiennon2022learning,ziegler2020finetuning}, story generation \citep{ziegler2020finetuning}, and adherence to instructions \citep{ouyang2022training,ramamurthy2023is}. The typical approach involves first learning a neural network-based reward model that fits the dataset of preference comparisons (often modeled using the Bradley-Terry framework \citep{bradley1952rankanalysis}), followed by optimizing the language model with reinforcement learning, such as REINFORCE \citep{williams1992reinforce}, proximal policy optimization (PPO; \cite{schulman2017proximal}), or related methods \citep{ramamurthy2023is}. Parallel efforts make use of instruction-tuned LLMs, further trained with human feedback, to generate new synthetic datasets reflecting specific attributes like safety or non-harmfulness \citep{bai2022constitutional}, relying mainly on loosely specified text-based rubrics provided by humans for supervision. Collectively, these strategies bridge two research directions: one focusing on the application of reinforcement learning to language modeling for a range of tasks~\citep{Ranzato2015SequenceLT,paulus2018a,wu2018learning}, and the other centered on systematic approaches to learning from human preferences \citep{christiano2017deep,kupcsik2018learning}. While exploiting relative human preference data offers considerable advantages, fine-tuning LLMs with reinforcement learning remains logistically difficult; the current study introduces a theoretically grounded method for directly optimizing relative preferences without resorting to RL.

Beyond linguistic applications, the problem of policy learning from preferences has been explored within both the bandit and reinforcement learning paradigms, leading to a variety of proposed methodologies. In particular, contextual bandit algorithms that utilize preference or ranking information, as opposed to explicit reward signals, are classified as contextual dueling bandits (CDB; \cite{yue2012karmed,dudik2015contextual}). When absolute rewards are unavailable, theoretical frameworks for CDBs rely on the concept of a \textit{von Neumann winner}—that is, a policy whose expected head-to-head win rate against every alternative policy meets or exceeds 50\% \citep{dudik2015contextual}. Notably, CDB methods assume that preference feedback is provided in an online fashion, whereas preference-based learning from human feedback frequently occurs with a predetermined batch of offline, annotated action pairs \citep{yan2022human}. Analogously, \textit{preference-based RL} (PbRL) utilizes binary preference information derived from an \textit{unknown} latent ‘scoring’ function rather than explicit reward values \citep{BusaFekete2014,ruiz2023dueling}. While several algorithms have been developed for PbRL—including approaches capable of incorporating off-policy preference data—the standard pipeline typically involves two steps: explicit estimation of the hidden scoring (reward) function followed by policy optimization with respect to this surrogate reward \citep{jain2013learning,BusaFekete2014,christiano2017deep,sadigh2017active,kupcsik2018learning}. In contrast, our work introduces a unified approach wherein a policy is directly trained to align with expressed preferences, foregoing intermediate reward model construction.

\section{Preliminaries}\label{section:prelims}

We summarize the RLHF framework as presented in \citeauthor{ziegler2020finetuning}, as well as in subsequent works \citep{stiennon2022learning, bai2022training, ouyang2022training}. This framework generally comprises three main components: (1) supervised fine-tuning (SFT); (2) collection of preferences and construction of a reward model; and (3) reinforcement learning-based policy optimization.

\textbf{SFT}: The initial stage of RLHF entails refining a large, pre-trained language model through supervised learning, using high-quality annotated datasets tailored to downstream tasks such as dialogue or summarization. This process yields a model denoted by $\pi^\text{SFT}$.

\textbf{Reward Modelling Phase}: The next step utilizes the SFT model, generating answer pairs $(y_1, y_2) \sim \pi^\text{SFT}(y \mid x)$ for given prompts $x$. These answer pairs are evaluated by human annotators, who indicate a preference—represented as $y_w \succ y_l \mid x$—with $y_w$ and $y_l$ denoting the more and less preferred responses, respectively. The assumption is that such preferences reflect an underlying but unobserved reward function $r^*(y, x)$. Several probabilistic models have been adopted for representing human preferences, with the Bradley-Terry (BT) model \cite{bradley1952rankanalysis} being especially common; broader Plackett-Luce ranking frameworks \citep{plackett1975analysis, luce2012individual} are also suitable when data includes multiple ranked completions. In the BT formulation, the human preference likelihood $p^*$ can be written as:

\begin{equation}\label{eq:bradley-terry}
    p^*(y_1\succ y_2 \mid x)=\frac{\exp\left(r^*(x, y_1)\right)}{\exp\left(r^*(x, y_1)\right) + \exp\left(r^*(x, y_2)\right)}.
\end{equation}

Given a fixed dataset of pairwise comparisons $\mathcal{D} = \bigl\{x^{(i)}, y_w^{(i)}, y_l^{(i)}\bigr\}_{i=1}^N$ sampled according to $p^*$, we can introduce a parameterized reward function $r_{\phi}(x, y)$ and learn its parameters through maximum likelihood estimation. Treating this as a binary classification problem leads to the following negative log-likelihood objective:

\begin{equation}\label{eq:reward_model}
    \mathcal{L}_R(r_{\phi}, \mathcal{D}) = -\mathbb{E}_{(x, y_w, y_l)\sim \mathcal{D}}\bigl[\log \sigma(r_{\phi}(x, y_w)- r_{\phi}(x, y_l))\bigr]
\end{equation}

where $\sigma$ denotes the logistic (sigmoid) function. Within language modeling frameworks, the function $r_{\phi}(x, y)$ is frequently initialized from the supervised fine-tuning (SFT) model $\pi^\text{SFT}(y \mid x)$, supplementing it with a final linear layer that outputs a single scalar reward for each input-output pair \cite{ziegler2020finetuning}. To minimize the variance of the reward estimates, it is customary to normalize the rewards such that $\mathbb{E}_{x,y\sim \mathcal{D}}\left[r_\phi(x, y)\right] = 0$ across all $x$, as previously proposed.

\textbf{RL Fine-Tuning Phase}: In the reinforcement learning phase, this reward function serves as a training signal for further language model refinement. In line with established methodologies~\citep{jaques2017sequence, jaques2020human}, the optimization objective is formulated as

\begin{equation}\label{eq:RL}
\max_{\pi_{\theta}}  \mathbb{E}_{x\sim \mathcal{D}, y\sim \pi_{\theta}(y \mid x)}\bigl[r_{\phi}(x, y)\bigr] - \beta\mathbb{D}_{\textrm{KL}}\bigl[\pi_{\theta}(y\mid x)\mid \mid \pi_\text{ref}(y\mid x)\bigr],
\end{equation}

where the scalar $\beta$ regulates divergence from a reference policy $\pi_\text{ref}$, typically the original SFT model $\pi^\text{SFT}$. In practice, the policy $\pi_\theta$ is also initialized with the weights of $\pi^\text{SFT}$. The KL constraint plays a crucial role by anchoring the learned policy near the reference distribution, thus ensuring the reliability of the reward model, preserving generation diversity, and preventing collapse onto a limited set of high-reward responses. Since language is inherently discrete, direct optimization of this objective is not tractable through standard gradient-based methods; instead, reinforcement learning is commonly used. The prevailing method \citep{ziegler2020finetuning, stiennon2022learning, bai2022training, ouyang2022training} reformulates the reward as ${r(x, y) = r_{\phi}(x, y) -\beta (\log \pi_{\theta}(y\mid x) - \log \pi_\text{ref}(y\mid x))}$, and optimizes it using PPO \cite{schulman2017proximal}. 

\section{Direct Preference Optimization}\label{sec:DPO}

Considering the substantial computational overhead of reinforcement learning in large-scale applications such as language model fine-tuning, we seek to introduce a streamlined policy optimization technique that utilizes preference data directly. In contrast to the conventional RLHF paradigm, which first trains an explicit reward function and then employs RL for policy improvement, our method adopts a specific parameterization of the reward such that the optimal policy can be derived analytically, circumventing the need for a separate RL optimization procedure. As we elaborate below, the principal innovation is to apply an explicit analytical mapping from reward functions to their associated optimal policies. This transformation enables conversion of reward-based objectives into policy-based objectives via a change of variables, eliminating the need to fit a standalone reward model, while retaining compatibility with well-established preference models like Bradley-T

\textbf{Derivation of the DPO Objective.} Our starting point is the RL objective with a general reward $r$, as formulated in Eq.~\ref{eq:RL}, consistent with earlier studies. Building upon foundational results from previous works~\citep{peters2007reinforcement, peng2019advantage, korbak2022reinforcement, go2023aligning}, one can readily verify that the solution optimizing the KL-regularized expected reward (as stated in Eq.~\ref{eq:RL}) is given by:

\begin{equation}\label{eq:op_policy}
    \pi_r(y\mid x) = \frac{1}{Z(x)}\pi_\text{ref}(y\mid x)\exp\left(\frac{1}{\beta}r(x, y)\right),
\end{equation}

where $Z(x) =\sum_{y}\pi_\text{ref}(y\mid x)\exp\left(\frac{1}{\beta}r(x, y)\right)$ denotes the partition function. For an explicit step-by-step derivation, refer to Appendix \ref{app:derivation1}. However, even when employing the MLE-estimated reward $r_\phi$ as an approximation of the true reward $r^*$, practical computation of $Z(x)$ remains computationally intensive~\citep{korbak2022reinforcement, go2023aligning}, limiting the usability of this parameterization. To address this, we can reformulate Eq.~\ref{eq:op_policy}, expressing $r$ in terms of the corresponding optimal policy $\pi_r$, the reference $\pi_\text{ref}$, and $Z(\cdot)$. By taking logarithms on both sides of Eq.~\ref{eq:op_policy} and manipulating terms, we arrive at:

\begin{equation}\label{eq:main_eq}
    r(x,y) =\beta \log \frac{\pi_r(y\mid x)}{\pi_\text{ref}(y\mid x)} + \beta \log Z(x).
\end{equation}

This transformation is applicable for both the ground-truth reward $r^*$ and its associated optimal policy $\pi^*$. Notably, since the Bradley-Terry model requires only the difference between two rewards, specifically ${p^*(y_1 \succ y_2 \mid x) = \sigma(r^*(x, y_1) - r^*(x, y_2))}$, substituting the expression from Eq.~\ref{eq:main_eq} into the preference model of Eq.~\ref{eq:bradley-terry} ensures the partition function $Z(x)$ cancels out. Consequently, the resulting probability of human preference can be solely represented in terms of $\pi^*$ and $\pi_\text{ref}$. Therefore, under the Bradley-Terry model, the optimal RLHF policy $\pi^*$ meets the preference condition:

\begin{equation}\label{eq:objective}
    p^*(y_1\succ y_2 \mid x)=\frac{1}{1 + \exp\left(\beta \log \frac{\pi^*(y_2\mid x)}{\pi_\text{ref}(y_2\mid x)} - \beta \log \frac{\pi^*(y_1\mid x)}{\pi_\text{ref}(y_1\mid x)}\right)}
\end{equation}

The full derivation can be found in Appendix~\ref{app:derivation2}. While this construction leverages the Bradley-Terry framework, the approach can be generalized, and analogous derivations for Plackett-Luce models~\citep{plackett1975analysis, luce2012individual} are detailed in Appendix~\ref{app:plackett_luce_models}.

This reformulation now expresses the likelihood of preference data in terms of the optimal policy, circumventing the need for the reward parameterization. It thus enables us to define a maximum likelihood estimation (MLE) criterion directly for a parameterized policy $\pi_\theta$. Mirroring the structure of the reward modeling objective (see Eq.~\ref{eq:reward_model}), the policy learning objective takes the form:

\begin{equation}\label{eq:optimum_model}
    \mathcal{L}_\text{DPO}(\pi_{\theta}; \pi_\text{ref}) = -\mathbb{E}_{(x, y_w, y_l)\sim \mathcal{D}}\left[\log \sigma \left(\beta \log \frac{\pi_{\theta}(y_w\mid x)}{\pi_\text{ref}(y_w\mid x

In this formulation, we construct an implicit reward function through an alternative parameterization, such that the policy attaining maximal performance is precisely $\pi_\theta$. Additionally, since this approach can be interpreted as fitting a reparameterized Bradley-Terry model, it inherits advantageous theoretical features, including consistency under suitable conditions on the distribution of preference data \cite{bong2022generalized}. We further elaborate on the theoretical aspects of DPO, drawing comparisons to related literature, in Section~\ref{sec:theory}.

\textbf{How does the DPO update operate?} To develop an intuition for DPO, it is instructive to examine the gradient of its loss function, $\mathcal{L}_\text{DPO}$. The derivative with respect to $\theta$ is expressed as:

\begin{multline*}\label{eq:gradient}
    \nabla_\theta \mathcal{L}_\text{DPO}(\pi_\theta;\pi_\text{ref}) = \\ -\beta\mathbb{E}_{(x, y_w, y_l) \sim \mathcal{D}} \bigg[\underbrace{\sigma(\hat{r}_\theta(x, y_l) - \hat{r}_\theta (x, y_w))}_\text{places higher emphasis when reward estimate is inaccurate}\bigg[\underbrace{\nabla_\theta\log \pi(y_w \mid x)}_\text{promotes higher probability for $y_w$} - \underbrace{\nabla_\theta\log\pi(y_l \mid x)}_\text{penalizes probability for $y_l$}\bigg]\bigg],
\end{multline*}

Here, $\hat{r}_\theta(x, y) = \beta \log \frac{\pi_\theta(y \mid x)}{\pi_\text{ref}(y \mid x)}$ represents the reward function implicitly induced by the language model $\pi_\theta$ relative to the reference model $\pi_\text{ref}$ (with further details provided in Section~\ref{sec:theory}). The gradient of $\mathcal{L}_\text{DPO}$ works by enhancing the probability of the preferred responses $y_w$ while suppressing that of the less preferred alternatives $y_l$. Notably, each instance is weighted according to the extent to which the implicit reward function $\hat{r}_\theta$ assigns a higher score to $y_l$ over $y_w$, modulated by $\beta$—essentially reflecting how erroneously the reward model orders the completions, and reflecting the influence of the KL regularization. Empirical findings indicate that this adaptive weighting plays a critical role; omitting it, as in a straightforward variant lacking the weighting term, may result in degeneration of the language model’s outputs (see Appendix Table~\ref{tab:unlikelihood_generations}).

\textbf{DPO overview.} The standard DPO workflow includes: 1) For each prompt $x$, draw sample completions $y_1, y_2 \sim \pi_\text{ref}(\cdot \mid x)$, assign human preference labels, and assemble the offline preference dataset $\mathcal{D} = \{x^{(i)}, y_w^{(i)}, y_l^{(i)}\}_{i=1}^N$; and 2) train the language model $\pi_\theta$ by minimizing $\mathcal{L}_\text{DPO}$ using the given $\pi_\text{ref}$, dataset $\mathcal{D}$, and the target $\beta$. 
In real-world scenarios, one typically prefers to leverage existing, publicly released preference datasets rather than collecting new samples and preferences. As these datasets have been generated by sampling from $\pi^\text{SFT}$, we set $\pi_\text{ref} = \pi^\text{SFT}$ whenever possible. If $\pi^\text{SFT}$ is not accessible, then $\pi_\text{ref}$ is instead initialized by maximizing the likelihood over preferred completions ${(x, y_w)}$, namely, by solving ${\pi_\text{ref} = \argmax_{\pi}\mathbb{E}_{x, y_w \sim \mathcal{D}}\left[\log \pi(y_w \mid x)\right]}$. This method addresses the potential distribution mismatch between the (unavailable) actual reference distribution and the operational $\pi_\text{ref}$ used within DPO. Appendix~\ref{app:implementation} provides additional specifics concerning implementation details and hyperparameter settings.

\section{Theoretical Analysis of DPO} Here, we deepen our interpretation of the DPO approach, supply theoretical justification, and connect the benefits of DPO to shortcomings seen in RLHF actor-critic algorithms (e.g., PPO~\cite{schulman2017proximal}).

\label{sec:theory}

\subsection{Your Language Model Is Secretly a Reward Model} DPO can forgo explicit reward fitting and policy optimization via RL, instead learning through a unified maximum likelihood framework. Observe that the optimization objective in Eq. \ref{eq:main_eq} corresponds to the Bradley-Terry model, where rewards are parameterized by $r^*(x, y) = \beta \log\frac{\pi^*_\theta(y \mid x)}{\pi_\text{ref}(y \mid x)}$. The process directly trains the model $\pi_\theta$ analogously to reward model training as in Eq. \ref{eq:reward_model}, with an appropriate change of variables. The following discussion develops the theoretical basis for this reparameterization, demonstrating that it places no restrictions on the set of learnable reward models and permits perfect recovery of the optimal policy. We begin by formalizing an equivalence relationship for reward functions.

\begin{definition}
Two reward functions $r(x, y)$ and $r'(x, y)$ are considered equivalent if and only if there exists a function $f$ such that ${r(x, y)-r'(x, y) = f(x)}$.
\end{definition}

It can be readily verified that this notion of equivalence defines an equivalence relation, thereby dividing the set of reward functions into distinct equivalence classes. We can formalize the following lemmas:

\begin{lemma}\label{lemma:same_prefrence} 
Within the Plackett-Luce preference structure, including the Bradley-Terry model as a special case, reward functions that belong to the same equivalence class will induce identical preference distributions.
\end{lemma}

\begin{lemma}\label{lemma:same_policy}
    Any pair of reward functions belonging to a single equivalence class will yield the same optimal policy when solving the corresponding constrained RL problem.
\end{lemma}

The arguments establishing these lemmas are straightforward and can be found in Appendix~\ref{app:lemma1}. The first lemma captures a well-documented ambiguity associated with the Plackett-Luce models \cite{plackett1975analysis}. Owing to this ambiguity, it is often necessary to introduce extra identifiability conditions in order to secure any guarantees for the maximum likelihood estimators obtained from Eq.~\ref{eq:reward_model} \cite{bong2022generalized}. According to the second lemma, all reward functions from a single equivalence class lead to the same optimal policy. Consequently, for the purpose of optimizing our final objective, it suffices to recover any member of the equivalence class containing the optimal reward. The result below, whose proof appears in Appendix~\ref{app:thm1}, formalizes this notion:

\begin{theorem}\label{thm:main}
    Assuming mild technical conditions, every equivalence class of reward functions consistent with the Plackett-Luce (and, specifically, Bradley-Terry) model can be expressed through the following reparameterization: ${r(x, y) = \beta \log \frac{\pi(y\mid x)}{\pi_\text{ref}(y\mid x)}}$ for an appropriate model $\pi(y\mid x)$ and some reference model $\pi_\text{ref}(y \mid x)$.
\end{theorem}

\begin{sproof}
    Take any reward function $r(x, y)$ that specifies an associated optimal model $\pi_r(y \mid x)$ as defined in Eq.~\ref{eq:op_policy}. We show that a reward function belonging to the same equivalence class as $r$ can be formulated using the above reparameterization. Define the projection operator $f$ by
\begin{equation}
    f(r; \pi_\text{ref}, \beta)(x, y) = r(x, y) - \beta\log\sum_{y}\pi_\text{ref}(y\mid x)\exp\left(\frac{1}{\beta}r(x, y)\right)
\end{equation}
This operator $f$ essentially normalizes the reward by subtracting the logarithm of the partition function with respect to $\pi_r$. Because the normalization depends solely on the prefix $x$, the function $f(r; \pi_\text{ref}, \beta)(x, y)$ lies in the same equivalence class as $r(x, y)$. If we substitute $r$ on the right-hand side of Eq.~\ref{eq:main_eq}, which is valid for any reward function, we obtain $f(r; \pi_\text{ref}, \beta)(x, y) = \beta \log \frac{\pi_r(y\mid x)}{\pi_\text{ref}(y\mid x)}$. Thus, $f$ produces a representative of the equivalence class of $r$ with the desired parametric form, and this reparameterization introduces no loss of generality for our reward model.
\end{sproof}

An alternative interpretation of Theorem~\ref{thm:main} is that it precisely identifies the particular reward function, within each equivalence class, that is chosen by the DPO reparameterization. This corresponds to the reward function that fulfills the following condition:

\begin{equation}\label{eq:lag_p}
     \sum_{y}\underbrace{\pi_\text{ref}(y\mid x)\exp\left(\frac{1}{\beta}r(x, y)\right)}_{=\pi(y\mid x)\text{, using Thm.~\ref{thm:main} reparam.}} = 1,
\end{equation}

In other words, $\pi(y\mid x)$ represents a legitimate probability distribution, meaning it assigns positive values that sum to 1. Furthermore, as seen from Eq.~\ref{eq:op_policy}, Eq.~\ref{eq:lag_p} describes the normalization factor (partition function) for the optimal policy associated with the reward function $r(x, y)$. The crucial insight provided by the DPO algorithm is that, by imposing certain constraints on the Plackett-Luce (and in particular, Bradley-Terry) preference model family—which is inherently under-constrained—we maintain the expressive capacity of possible reward models, while also ensuring that the optimal policy expressed in Eq.~\ref{eq:op_policy} remains analytically computable for any prompt $x$.

\subsection{Instability of Actor-Critic Algorithms}
Our framework also enables the identification of sources of instability in commonly employed actor-critic methods for RLHF, such as PPO. Adopting the RLHF methodology, we focus on the RL fine-tuning phase detailed in Section~\ref{section:prelims}. This approach allows us to relate to the control-as-inference paradigm \cite{levine2018reinforcement} for the constrained reinforcement learning setup introduced in \ref{eq:RL}. Given a parameterized model $\pi_{\theta}(y\mid x)$, the objective is to minimize $\mathbb{D}_{\text{KL}}[\pi_{\theta}(y|x) \mid \mid \pi^*(y\mid x)]$, where $\pi^*$ denotes the optimal policy from Eq.~\ref{eq:optimum_model} determined by the reward $r_{\phi}(y, x)$. Through some mathematical manipulations, this yields the following objective function:

\begin{equation}\label{eq:AC}
    \max_{\pi_{\theta}}\mathbb{E}_{\pi_{\theta}(y\mid x)}\bigg[\underbrace{r_{\phi}(x, y) -\beta\log\sum_{y}\pi_\text{ref}(y\mid x)\exp\left(\frac{1}{\beta}r_{\phi}(x, y)\right)}_{f(r_{\phi}, \pi_\text{ref}, \beta)} - \underbrace{\beta\log\frac{\pi_{\theta}(y\mid x)}{\pi_\text{ref}(y\mid x)}}_{\text{KL}}\bigg]
\end{equation}

The objective addressed here aligns with that of previous studies 
\citep{ziegler2020finetuning, stiennon2022learning, bai2022training, ouyang2022training}, which utilize the DPO-equivalent reward within the $r_{\phi}$ reward class. In this context, the normalization component in $f(r_{\phi}, \pi_\text{ref}, \beta)$ may be viewed as the soft value function associated with the reference policy $\pi_\text{ref}$. Although this normalization does not alter the location of the optimal solution, omitting it can substantially increase the variance of the policy gradient estimator, thereby introducing instability during training. To manage this, one might employ a learned value function to approximate the normalization, but optimizing such a function can itself pose significant challenges. A common alternative in prior research involves using the reward associated with a human completion as a normalization baseline, effectively yielding a one-sample Monte Carlo estimate for the normalization term. The DPO reparameterization, however, gives rise to a reward formulation that eliminates the need for explicit baselines.

\section{Experiments}
This section presents an empirical investigation of DPO, focusing on its capacity to directly train policies from preference data. We begin with controlled experiments in text generation to address the question: how effectively does DPO manage the trade-off between reward maximization and maintaining low KL-divergence from the reference policy, in comparison with established preference learning methods like PPO? Subsequently, we assess the scalability of DPO on larger model architectures and more challenging RLHF applications such as summarization and dialogue systems. Remarkably, our results indicate that, even with minimal hyperparameter adjustment, DPO typically matches or surpasses the performance of strong benchmarks—specifically PPO-based RLHF—as well as surpassing strategies that select the highest reward trajectory among $N$ samples generated under a learned reward model. The following subsections detail the experimental methodology, with supplementary information available in Appendix~\ref{app:exp_details}.

\textbf{Tasks.} We investigate three distinct open-ended text generation tasks in our experiments. Across all tasks, the models are trained to optimize a policy based on a set of preference data, $\mathcal{D}=\bigl\{x^{(i)}, y_w^{(i)}, y_l^{(i)}\bigr\}_{i=1}^N$. For the \textbf{controlled sentiment generation} task, the input $x$ consists of an initial segment from a movie review in the IMDb dataset \cite{maas-EtAl:2011:ACL-HLT2011}, and the objective is to generate $y$ that exhibits positive sentiment. To enable rigorous evaluation, we automatically construct preference pairs between generations using a sentiment classifier such that $p(\text{positive}\mid x,y_w)>p(\text{positive}\mid x,y_l)$. The SFT baseline involves fine-tuning GPT-2-large on the training portion of the IMDb reviews dataset until convergence (additional details are available in App~\ref{app:sentiment_details}). For the \textbf{summarization} task, each input $x$ is a Reddit forum post, and the policy must summarize its main points in $y$. Following established methodology, we leverage the Reddit TL;DR summarization dataset \citep{volske-etal-2017-tl} alongside human preferences previously collected by \citeauthor{stiennon2022learning}. Our SFT baseline is a model fine-tuned on summaries written by humans for Reddit posts\footnote{\url{https://huggingface.co/CarperAI/openai_summarize_tldr_sft}}, implemented using the TRLX \citep{leandro_von_werra_2023_7790115} library to support RLHF. The human preference annotations correspond to samples generated by another SFT model, similar in training to our own, as in \citeauthor{stiennon2022learning}. For the \textbf{single-turn dialogue} task, each prompt $x$ is a human-issued query, ranging from scientific inquiries to personal advice. The goal is to generate an informative and engaging reply $y$. To this end, we use the Anthropic Helpful and Harmless dialogue dataset \citep{bai2022training}, which comprises 170,000 human–assistant conversation transcripts, each concluding with two model-generated candidate responses and an associated human preference. Since no SFT checkpoint is provided for this task, we build our SFT model by fine-tuning an existing language model solely on the subset of preferred responses.

\textbf{Evaluation.} Our experimental framework incorporates two distinct evaluation methodologies. To assess each algorithm's capability in optimizing the constrained reward maximization objective, we perform a controlled sentiment generation analysis in which we chart the Pareto frontier between obtained reward and KL-divergence relative to the reference policy; this is feasible since the underlying reward function (provided by a sentiment classifier) is fully observable. Conversely, in more realistic settings where the true reward function is inaccessible, we measure algorithmic performance by comparing their \textit{win rate} to a baseline policy. This is accomplished via GPT-4, serving as a stand-in for human assessors when judging the quality of summaries and the helpfulness of responses in summarization and single-turn dialogue tasks, respectively. For summarization experiments, the reference summaries within the test corpus establish the baseline, whereas in dialogue evaluations, the baseline comprises the test set’s preferred responses. Although literature indicates that LMs may surpass standard automated metrics in evaluation quality \citep{Chen2023ExploringTU}, we augment our results with a human subject study, as detailed in Sec.~\ref{sec:human-judgments}, to validate the use of GPT-4 for evaluative purposes. Our analysis reveals that GPT-4’s decisions show strong concordance with human raters, with human-GPT-4 agreement on par with, or exceeding, inter-human annotator agreement.

\textbf{Methods.} Alongside DPO, we compare a variety of established techniques for training language models in accordance with human preferences. For the summarization task, we employ zero-shot prompting using \textbf{GPT-J} \citep{gpt-j}, while the dialogue task utilizes 2-shot prompting with \textbf{Pythia-2.8B} \citep{biderman2023pythia}. Further, we assess the \textbf{SFT} model and the \textbf{Preferred-FT} variant, the latter involving supervised fine-tuning on preferred completions $y_w$ sourced from either the SFT model (in controlled sentiment and summarization) or from a standard LM (for single-turn dialogue). The \textbf{Unlikelihood} approach~\citep{welleck2019neural} offers another pseudo-supervised strategy, where the model is explicitly trained to increase the probability of $y_w$ and \textit{decrease} that of $y_l$; an adjustable parameter $\alpha \in [0,1]$ weights the unlikelihood component. Our benchmarks also feature \textbf{PPO} \citep{schulman2017proximal}, relying on a reward model trained from preference annotations, and \textbf{PPO-GT}, an oracle variant learning from the ground-truth reward in the controlled sentiment scenario. Within sentiment experiments, we implement two versions of PPO-GT: an off-the-shelf release \cite{leandro_von_werra_2023_7790115}, and an enhanced version featuring reward normalization and optimized hyperparameters to further elevate results (these improvements also inform our standard PPO runs with learned rewards). Lastly, the \textbf{Best of $N$} approach serves as a baseline wherein $N$ completions are sampled from the SFT (or Preferred-FT for dialogue), and the response with the highest reward—according to the learned reward model—is selected. While this method typically yields strong outcomes and avoids entangling the reward and PPO optimization processes, it is computationally prohibitive for even modest values of $N$, given the need to produce $N$ completions per evaluation instance.

\subsection{How well can DPO optimize the RLHF objective?}

The objective in standard RLHF algorithms involves maximizing rewards while enforcing a KL divergence constraint to limit the policy’s divergence from the reference policy, essentially mediating between reward exploitation and policy conservativeness. Therefore, algorithm comparisons should consider both the magnitude of reward obtained and the extent of KL divergence; achieving a marginally higher reward at the cost of significantly increased KL may not be preferable. Figure~\ref{fig:frontier-tldr-main} presents the trade-off curve between reward and KL for different methods within the sentiment analysis context. For each method, we conduct multiple training runs, each with distinct policy conservativeness hyperparameters (PPO: target KL $\in\{3,6,9,12\}$; DPO: $\beta \in \{0.05,0.1,1,5\}$; unlikelihood: $\alpha\in\{0.05,0.1,0.5,1\}$; preferred-FT: different random seeds). Altogether, 22 configurations are assessed. After every 100 training steps, up to convergence, we evaluate policies over a collection of test prompts, calculating the mean reward according to the true reward metric and the mean sequence-level KL\footnote{Defined as the cumulative sum of the timestep-wise KL divergences.} relative to the reference policy, $\text{KL}\left(\pi\mid \mid \pi_\text{ref}\right)$. The results indicate that DPO achieves a frontier that is substantially more favorable—yielding superior rewards at comparably low KL values. This outcome is significant for several reasons. Although DPO and PPO share the same optimization objective, DPO demonstrates marked efficiency gains: its reward/KL curve consistently surpasses that of PPO. Furthermore, DPO outperforms PPO even when the latter has access to ground-truth rewards (PPO-GT), resulting in a strictly better reward-KL frontier.

\subsection{Can DPO scale to real preference datasets?}
\label{sec:dpo-real-datasets}
We proceed by assessing how well DPO performs when fine-tuned on tasks in summarization and single-turn dialogue using real-world preference data. In the context of summarization, standard automatic evaluation metrics like ROUGE often fail to align closely with human judgment~\citep{stiennon2022learning}, and earlier studies have shown that preference-driven fine-tuning of language models with PPO can yield more preferable summaries. For our experiments, we generate completions for various approaches on the test set of the TL;DR summarization dataset, then measure the average rate at which each method’s completions are favored over the provided reference completions. Sampling is performed across a temperature range from 0.0 to 1.0, and the corresponding win rates are illustrated in Figure~\ref{fig:frontier-tldr-main} (right). All techniques—including DPO, PPO, and Preferred-FT—utilize fine-tuning starting from the same GPT-J SFT model\footnote{\url{https://huggingface.co/CarperAI/openai_summarize_tldr_sft}}. Our results indicate that at a sampling temperature of 0.0, DPO achieves a win rate of roughly 61\%, surpassing PPO’s best result of about 57\% at the same temperature. Additionally, DPO attains a higher peak win rate than the best among the $N$ baseline models. It is worth mentioning that the $\beta$ hyperparameter in DPO was not extensively optimized, which suggests these results could underestimate DPO’s full potential. DPO also demonstrates greater robustness to changes in sampling temperature relative to PPO; notably, PPO’s performance diminishes to that of the original GPT-J model when the temperature increases. Meanwhile, Preferred-FT yields minimal improvement over the SFT model. Further, direct comparison between DPO and PPO using human evaluators is presented in Section~\ref{sec:human-judgments}, where DPO samples at temperature 0.25 are preferred in 58\% of cases compared to PPO samples at the same temperature.

For the single-turn dialogue scenario, we assess various approaches on the portion of the Anthropic HH test split \citep{bai2022training} involving a single human-assistant exchange. In the GPT-4-based assessments, win rates are determined by comparing the model outputs against the preferred completions in the test set. Due to the absence of a standard SFT baseline for this problem, we initialize with a pre-trained Pythia-2.8B model, train a reference model via Preferred-FT using the preferred responses—ensuring these responses are representative for the model’s distribution—and then proceed to train with DPO. Additional comparisons include taking the best output among 128 Preferred-FT completions (since our analyses show that the Best of $N$ method saturates at $N=128$ for this dataset; refer to Appendix Figure~\ref{fig:best-of-n}), as well as evaluating a two-shot prompted variant of the Pythia-2.8B base. Across optimal sampling temperatures for each technique, DPO consistently matches or exceeds the performance of alternatives. We further test an RLHF model trained with PPO on the Anthropic HH data \footnote{\url{https://huggingface.co/reciprocate/ppo_hh_pythia-6B}}, using resources from a prominent repository \footnote{\url{https://github.com/CarperAI/trlx/tree/main/examples/hh}}; however, suitable prompts and temperatures for this model fail to surpass the original Pythia-2.8B base model’s performance. Drawing from TL;DR results and recognizing that both approaches maximize the same reward signal, we treat the Best of 128 strategy as a practical stand-in for PPO-level performance. Summing up, DPO stands out as the only resource-efficient method delivering improvements over the preferred completions in Anthropic HH, while also matching or outperforming the computationally intensive Best of 128 baseline. Notably, Figure~\ref{fig:dialogue-main} illustrates that DPO attains optimal performance rapidly.

\subsection{Generalization to a new input distribution}

To more thoroughly assess how PPO and DPO perform under distributional changes, we take the policies trained on the Reddit TL;DR summarization task and evaluate them on a novel dataset: news articles from the test set of CNN/DailyMail \citep{nallapati-etal-2016-abstractive}, applying the optimal sampling temperatures established for TL;DR (0 and 0.25). The corresponding outcomes are displayed in Table~\ref{tab:ood}. The comparison metric is the GPT-4 win rate versus the reference summaries from the dataset, using the same GPT-4 (C) prompt as for Reddit TL;DR, substituting ``forum post'' with ``news article.'' On this out-of-distribution task, DPO maintains a clear performance advantage over the PPO policy. These findings suggest that DPO policies possess similar, if not superior, generalization capabilities as PPO policies, despite the fact that DPO does not leverage the supplementary unlabeled Reddit TL;DR prompts accessible to PPO.

\subsection{Validating GPT-4 judgments with human judgments}
\label{sec:human-judgments}
We implement a human evaluation to examine the validity of GPT-4's comparative assessments, focusing on the TL;DR summarization setting with two distinct GPT-4 prompting styles. The \textbf{GPT-4 (S)} (simple) prompt directs the model to judge which summary better captures the core content. In contrast, the \textbf{GPT-4 (C)} (concise) prompt explicitly instructs the model to consider conciseness, as we observed that GPT-4 with the \textbf{GPT-4 (S)} prompt often exhibits a bias toward longer, more repetitive summaries compared to human preferences. Full prompt texts are provided in Appendix~\ref{app:prompts}. We undertake three paired comparisons: the best-performing model (DPO, temp. 0.25), the lowest-performing (PPO, temp. 1.0), and an intermediate approach (SFT, temp. 0.25), each pitted against PPO using its best greedy decoding (i.e., optimal temperature). This selection is intended to sample a broad range of summary qualities; for DPO and PPO-1, multiple human assessments are collected, but limited rater availability prevents further duplication. Our results indicate that, across both prompts, GPT-4's preferences align with human evaluators about as frequently as humans agree with one another, supporting the use of GPT-4 as a stand-in for human judgment. Notably, the \textbf{GPT-4 (C)} prompt yields win rates that more closely mirror human outcomes, motivating its use as the main evaluation metric in Section~\ref{sec:dpo-real-datasets}. For expanded information about the human evaluation protocol, rater interface, and volunteer roster, please refer to Appendix~\ref{app:human-study}.

\section{Discussion}
Preference-based learning presents a robust and scalable approach for developing language models that are both proficient and aligned with user values. In this work, we proposed DPO, a straightforward methodology for optimizing language models using preference data, entirely circumventing reinforcement learning. Instead of adapting preference learning to fit within the traditional RL paradigm—and consequently relying on existing RL algorithms—DPO leverages a principled correspondence between language model policies and reward structures. This framework enables direct training of models to reflect human preferences through a basic cross-entropy loss, foregoing the need for reinforcement learning or compromising on model expressiveness. Remarkably, DPO achieves performance on par with, or superior to, leading RLHF methods such as PPO, with almost negligible hyperparameter adjustment. Consequently, DPO significantly lowers the complexity involved in preference-based language model training.

\textbf{Limitations \& Future Work.} Our findings prompt a range of pertinent directions for subsequent investigation. For instance, it remains to be clarified how the generalization of DPO-trained policies compares with models optimized against explicit reward signals, particularly in out-of-distribution settings. Preliminary evidence indicates comparable generalization capabilities relative to PPO-optimized policies, but this assertion warrants deeper exploration. Additional open questions include whether utilizing DPO for self-labeling allows for efficient exploitation of unlabeled prompt data. Further, it is important to understand the dynamics of reward over-optimization within the direct preference optimization context—such as whether the moderate performance drop observed in Figure~\ref{fig:dialogue-main}-right exemplifies this phenomenon. Moreover, while our experiments address models up to 6B parameters, scaling DPO for use with significantly larger, state-of-the-art architectures represents a promising avenue. On the evaluation front, we note that GPT-4-determined win rates are sensitive to prompt wording, suggesting future studies should aim to determine optimal methodologies for soliciting reliable judgments from automated evaluators. Finally, DPO holds substantial potential for broader applications, extending to the training of generative models across different domains beyond language modeling.